% !TeX root = main.tex
\section{Limitations and Discussion}

잔여 셰이핑 가치의 감소는 지침의 유해성과 동일하지 않다. 서브골이 성공적으로
달성되었거나 현재 정책이 해당 사건을 아직 수행하지 못하는 경우에도 예측값은 낮아질
수 있다. 또한 critic의 타깃에는 시간에 따라 변하는 정책과 서로 다른 지침 아래에서
수집된 궤적이 섞이므로, $\widehat V_F$를 지침 유지의 반사실 가치로 해석할 수 없다.
Episode 경계와 critic drift 역시 비율 및 검출 통계량에 영향을 줄 수 있다.

현재 grounding은 제한된 수의 사람이 정의한 predicate에 의존한다. 따라서 LLM이
생성한 복합 조건, 시간 범위, 새로운 전략이 실행 의미로 보존되지 않을 수 있다.
RSVP도 grounding 라이브러리에 존재하지 않는 서브골은 감시할 수 없다. 이는 LLM의
표현력을 활용하면서 계산 가능성과 검증 가능성을 유지해야 하는 더 근본적인 문제다.

또한 식~\eqref{eq:reward}의 셰이핑 보상은 일반적인 potential-based reward
형태로 구성되지 않았다~\cite{ng1999shaping}. 학습 후반에 $\lambda_t$를 줄이는 것은
보조 목적에 대한 의존을 완화할 수 있지만, 원래 환경에서의 최적 정책이 보존된다는
이론적 보장을 제공하지는 않는다.

마지막으로 adaptive refresh가 fixed refresh보다 항상 좋거나 호출 예산을 자동으로
절감한다는 보장은 없다. 조기 호출은 같은 $F_{\max}$의 고정 기준선보다 호출을 늘릴
수 있다. 따라서 성능과 비용을 함께 보고하고, 호출량을 맞춘 기준선과 비교해야 한다.
이러한 한계는 부가적인 주의사항이 아니라 본 방법의 성립 여부를 결정하는 검증 대상이다.

\section{Conclusion}

본 연구는 LLM 지침이 결합된 MARL에서 지침의 내용과 무관한 고정 갱신 주기의 한계를
다룬다. RSVP는 보유 지침이 강조하는 predicate의 잔여 셰이핑 가치를 예측하고 그
감소가 누적될 때 지침을 조기 갱신한다. 이 설계는 기존 grounding 신호를 재사용하여
별도의 LLM 평가 호출을 요구하지 않으며, 최대 갱신 간격을 fallback으로 유지한다.
향후 실험은 기반 LLM 지침의 실제 효용, critic의 시간순 예측력, 검출 후 갱신의 외부
성능 이득을 순차적으로 검증할 것이다. 이 세 연결고리가 확인될 때 RSVP는 고정 주기
LLM 지침 시스템을 내용 인지형 온디맨드 갱신으로 확장하는 근거를 갖게 된다.
